\documentclass{book}

\title{An Introduction to Tensor Calculus}
\newcommand{\booksubtitle}{A Latex Template for a Science Textbook}
\newcommand{\booklicense}{Creative Commons Zero 1.0 Universal}

\author{Author Name}
% Author subtitle could be a university or a geographical location, for example
\newcommand{\authorsubtitle}{City, Country}

% Create convenient commands \booktitle and \bookauthor
\makeatletter
\newcommand{\booktitle}{\@title}
\newcommand{\bookauthor}{\@author}
\makeatother

% This utf8 declaration is not needed for versions of latex > 2018 but may
% be helpful for older software. Eventually it may not be worth keeping.
\usepackage[utf8]{inputenc}
\usepackage{fix-cm}  % this package allows large \fontsize
\usepackage{tikz}    % this is for graphics. e.g. rectangle on title page
\usepackage{amsmath, amssymb} % Used by equations

% The following dimensions specify 4.75" X 7.5" content on 6 3/8" by 9 1/4"
% paper. The paper width and height can be tweaked as required and the content
% should size to fit within the margins accordingly.
%
% The (inside) bindingoffset should be larger for books with more pages. Some
% standard recommended sizes are .375in minimum up to 1in for 600+ page books.
% Sizes .75in and .875in are also recommended roughly at 150 and 400 pages.
\usepackage[bindingoffset=0.75in,
            left=.5in, right=.5in,
            top=.8125in, bottom=.9375in,
            paperwidth=6.375in, paperheight=9.25in]{geometry}
% Here is an alternative geometry for reading on letter size paper:
% \usepackage[margin=.75in, paperwidth=8.5in, paperheight=11in]{geometry}

\renewcommand{\contentsname}{Table of Contents} % default is {Contents}
\usepackage{makeidx}
\makeindex % Initialize an index so we can add entries with \index

\begin{document}
\frontmatter

% ---- Half Title Page ----
% current geometry will be restored after title page
\newgeometry{top=1.75in,bottom=.5in}
\begin{titlepage}
\begin{flushleft}

% Title
\textbf{\fontfamily{qcs}\fontsize{48}{54}\selectfont Science\\Textbook\\}

% Draw a line 4pt high
\par\noindent\rule{\textwidth}{4pt}\\

% Subtitle
% Shaded box from left to right. Text node is midway (centered).
\begin{tikzpicture}
\shade[bottom color=lightgray,top color=white]
    (0,0) rectangle (\textwidth, 1.5)
    node[midway] {\textbf{\large \textit{\booksubtitle}}};
\end{tikzpicture}

% Edition Number
\begin{flushright}
\Large First Edition
\end{flushright}

% \vspace{\fill}
\vspace{\fill}

\end{flushleft}
\end{titlepage}
\restoregeometry
% ---- End of Half Title Page ----


% No page numbers on the Frontispiece page
\thispagestyle{empty}


% ---- Title Page ----
% current geometry will be restored after title page
\newgeometry{top=1.75in,bottom=.5in}
\begin{titlepage}
\begin{flushleft}

% Title
\textbf{\fontfamily{qcs}\fontsize{48}{54}\selectfont Science\\Textbook\\}

% Draw a line 4pt high
\par\noindent\rule{\textwidth}{4pt}\\

% Shaded box from left to right with Subtitle
% The text node is midway (centered).
\begin{tikzpicture}
\shade[bottom color=lightgray,top color=white]
    (0,0) rectangle (\textwidth, 1.5)
    node[midway] {\textbf{\large \textit{\booksubtitle}}};
\end{tikzpicture}

% Edition Number
\begin{flushright}
\Large First Edition
\end{flushright}

\vspace{\fill}

% Author and Location
\textbf{\large \bookauthor}\\[3.5pt]
\textbf{\large \textit{\authorsubtitle}}

\vspace{\fill}

% Self Publishing Logo. Free to use: CC0 license.
% The source file is book.svg. If you change the svg, you must then convert
% it to pdf. There are many online and offline tools available to do that.
\begin{center}
%\includegraphics{booksvg.pdf}\\[4pt]
\fontfamily{lmtt}\small{Self Publishers Worldwide\\
Seattle San Francisco New York\\
London Paris Rome Beijing Barcelona}
\end{center}

\end{flushleft}
\end{titlepage}
\restoregeometry
% ---- End of Title Page ----

% Do not show page numbers on colophon page
\thispagestyle{empty}

\begin{flushleft}
\vspace*{\fill}
This book was typeset using \LaTeX{} software.\\
\vspace{\fill}
Copyright \textcopyright{} \the\year{}  \bookauthor\\
License: \booklicense
\end{flushleft}

% A title page resets the page # to 1, but the second title page
% was actually page 3. So add two to page counter.
\addtocounter{page}{2}

% The asterisk excludes chapter from the table of contents.
\chapter*{Preface}

\[
\left.
\begin{aligned}
  &\sqrt{{dX_{1}}^{2} + {dX_{2}}^{2} + {dX_{3}}^{2}} \\
  &\qquad
  \begin{aligned}
  &= \left(1 + \frac{\kappa}{8\pi} \int \frac{\sigma\, dV_{0}}{r}\right)
     \sqrt{{dx_{1}}^{2} + {dx_{2}}^{2} + {dx_{3}}^{2}}, \\
dT &= \left(1 - \frac{\kappa}{8\pi} \int \frac{\sigma\, dV_{0}}{r}\right) dl.
\end{aligned}
\end{aligned}
\right\}
\]

% Three-level Table of Contents
\setcounter{tocdepth}{3}
\tableofcontents

\mainmatter

\chapter{Prerequisites}

The theory of relativity is intimately connected
with the theory of space and time. I shall therefore begin with a brief
investigation of the origin of our ideas of space and time, although in
doing so I know that I introduce a controversial subject. The object of all
science, whether natural science or psychology, is to co-ordinate our
experiences and to bring them into a logical system. How are our customary
ideas of space and time related to the character of our experiences?

\section{First Principles}

The theory of relativity is intimately connected
with the theory of space and time. I shall therefore begin with a brief
investigation of the origin of our ideas of space and time, although in
doing so I know that I introduce a controversial subject. The object of all
science, whether natural science or psychology, is to co-ordinate our
experiences and to bring them into a logical system. How are our customary
ideas of space and time related to the character of our experiences?

\subsection{Examples}

The theory of relativity is intimately connected
with the theory of space and time. I shall therefore begin with a brief
investigation of the origin of our ideas of space and time, although in
doing so I know that I introduce a controversial subject. The object of all
science, whether natural science or psychology, is to co-ordinate our
experiences and to bring them into a logical system. How are our customary
ideas of space and time related to the character of our experiences?

\subsection{Excessive Elaborations}

The theory of relativity is intimately connected
with the theory of space and time. I shall therefore begin with a brief
investigation of the origin of our ideas of space and time, although in
doing so I know that I introduce a controversial subject. The object of all
science, whether natural science or psychology, is to co-ordinate our
experiences and to bring them into a logical system. How are our customary
ideas of space and time related to the character of our experiences?

\subsection{Long Winded Conclusion}

The theory of relativity is intimately connected
with the theory of space and time. I shall therefore begin with a brief
investigation of the origin of our ideas of space and time, although in
doing so I know that I introduce a controversial subject. The object of all
science, whether natural science or psychology, is to co-ordinate our
experiences and to bring them into a logical system. How are our customary
ideas of space and time related to the character of our experiences?

\subsection{Exercises}

The theory of relativity is intimately connected
with the theory of space and time. I shall therefore begin with a brief
investigation of the origin of our ideas of space and time, although in
doing so I know that I introduce a controversial subject. The object of all
science, whether natural science or psychology, is to co-ordinate our
experiences and to bring them into a logical system. How are our customary
ideas of space and time related to the character of our experiences?

\section{Theory vs Practice}

The theory of relativity is intimately connected
with the theory of space and time. I shall therefore begin with a brief
investigation of the origin of our ideas of space and time, although in
doing so I know that I introduce a controversial subject. The object of all
science, whether natural science or psychology, is to co-ordinate our
experiences and to bring them into a logical system. How are our customary
ideas of space and time related to the character of our experiences?

\subsection{Examples}

The theory of relativity is intimately connected
with the theory of space and time. I shall therefore begin with a brief
investigation of the origin of our ideas of space and time, although in
doing so I know that I introduce a controversial subject. The object of all
science, whether natural science or psychology, is to co-ordinate our
experiences and to bring them into a logical system. How are our customary
ideas of space and time related to the character of our experiences?

\subsection{Excessive Elaborations}

The theory of relativity is intimately connected
with the theory of space and time. I shall therefore begin with a brief
investigation of the origin of our ideas of space and time, although in
doing so I know that I introduce a controversial subject. The object of all
science, whether natural science or psychology, is to co-ordinate our
experiences and to bring them into a logical system. How are our customary
ideas of space and time related to the character of our experiences?

\subsection{Long Winded Conclusion}

The theory of relativity is intimately connected
with the theory of space and time. I shall therefore begin with a brief
investigation of the origin of our ideas of space and time, although in
doing so I know that I introduce a controversial subject. The object of all
science, whether natural science or psychology, is to co-ordinate our
experiences and to bring them into a logical system. How are our customary
ideas of space and time related to the character of our experiences?

\subsection{Exercises}

The theory of relativity is intimately connected
with the theory of space and time. I shall therefore begin with a brief
investigation of the origin of our ideas of space and time, although in
doing so I know that I introduce a controversial subject. The object of all
science, whether natural science or psychology, is to co-ordinate our
experiences and to bring them into a logical system. How are our customary
ideas of space and time related to the character of our experiences?

\foreach \x in {A,B,C,D,E,F,G,H,I,J,K,L,M,N,O,P,Q,R,S,T,U,V,W,X,Y,Z}
    {\index{\x 1}\index{\x 2}}

\chapter{Tensors}



\chapter{Metric Tensor}

Consider a vector that is represented in the Cartesian coordinate system and ***

You can use the Pythagoras' Theorem \(a^2+b^2=c^2\) to find its length

\begin{align*}
  ||\vec{v}||^2
  &= \left(v^1\right)^2+\left(v^2\right)^2 \\
  &= \left(3\right)^2+\left(4\right)^2 \\
  &= 25
\end{align*}
\begin{align*}
  \therefore ||\vec{v}||=\sqrt{25}=5
\end{align*}


But what if the vector is measured in this coordinate system? How would you measure its length? Let's try using Pythagoras' theorem and see if we get the same answer.

 \begin{align*}
||\vec{v}||^2
&= \left(\tilde{v}^1\right)^2+\left(\tilde{v}^2\right)^2 \\
&= \left(\frac{7}{5}\right)^2+\left(\frac{13}{5}\right)^2 \\
&= \frac{218}{25}
\end{align*}

\begin{align*}
\therefore ||\vec{v}|| = \sqrt{\frac{218}{25}} \approx 2.95?
\end{align*}

We got a different answer this time because Pythagoras' theorem only holds for the Cartesian coordinate system. In order to find the length of a vector in another coordinate system, recall that:

 \begin{align*}
   ||\vec{v}||^2
   = \vec{v}\cdot\vec{v}
 \end{align*}

 If we represent the vector in terms of the Cartesian coordinate system, we get the Pythagoras' Theorem back.

\begin{align*}
  \vec{v}\cdot\vec{v}
  &= (v^1\vec{e}_1+v^2\vec{e}_2) \cdot (v^1\vec{e}_1+v^2\vec{e}_2) \\
  &= v^1v^1(\vec{e}_1\cdot \vec{e}_1)+v^1v^2(\vec{e}_1\cdot \vec{e}_2)+v^2v^1(\vec{e}_2\cdot \vec{e}_1)+v^2v^2(\vec{e}_2\cdot \vec{e}_2) \\
  &= (v^1)^2(\vec{e}_1\cdot \vec{e}_1)+2v^1v^2(\vec{e}_1\cdot \vec{e}_2)+(v^2)^2(\vec{e}_2\cdot \vec{e}_2) \\
  &= (v^1)^2+(v^2)^2 \tag{Because \(\vec{e}_i\cdot\vec{e}_j = \delta_{ij}\)}
 \end{align*}

But if we represent the vector in terms of the basis vectors \(\tilde{e_1}\) and \(\tilde{e_2}\) we get:

\begin{align}
  \vec{v}\cdot\vec{v}
  &= (\tilde{v}^1\tilde{e}_1+\tilde{v}^2\tilde{e}_2) \cdot (\tilde{v}^1\tilde{e}_1+\tilde{v}^2\tilde{e}_2) \nonumber\\
  &= \tilde{v}^1\tilde{v}^1(\tilde{e}_1\cdot \tilde{e}_1)+\tilde{v}^1\tilde{v}^2(\tilde{e}_1\cdot \tilde{e}_2)+\tilde{v}^2\tilde{v}^1(\tilde{e}_2\cdot \tilde{e}_1)+\tilde{v}^2\tilde{v}^2(\tilde{e}_2\cdot \tilde{e}_2) \nonumber \\
  &= (\tilde{v}^1)^2(\tilde{e}_1\cdot \tilde{e}_1)+2\tilde{v}^1\tilde{v}^2(\tilde{e}_1\cdot \tilde{e}_2)+(\tilde{v}^2)^2(\tilde{e}_2\cdot \tilde{e}_2)
\end{align}

To find \((\tilde{e}_1\cdot \tilde{e}_1)\), \((\tilde{e}_1\cdot \tilde{e}_2\)),\((\tilde{e}_2\cdot \tilde{e}_2)\) we represent them in Cartesian coordinates first

\begin{align*}
  \tilde{e}_1\cdot\tilde{e}_1
  &=(2\vec{e}_1+1\vec{e}_2)\cdot(2\vec{e}_1+1\vec{e}_2)\\
  &=2^2(\vec{e}_1\cdot\vec{e}_1)+2(2)(1)(\vec{e}_1\cdot\vec{e}_2)+1^2(\vec{e}_2\cdot\vec{e}_2)\\
  &=2^2+1=5
\end{align*}

\begin{align*}
  \tilde{e}_1\cdot\tilde{e}_1
  &=(2\vec{e}_1+1\vec{e}_2)\cdot(2\vec{e}_1+1\vec{e}_2)\\
  &=2^2(\vec{e}_1\cdot\vec{e}_1)+2(2)(1)(\vec{e}_1\cdot\vec{e}_2)+1^2(\vec{e}_2\cdot\vec{e}_2)\\
  &=2^2+1=5
\end{align*}

\begin{align*}
  \tilde{e}_1\cdot\tilde{e}_1
  &=(2\vec{e}_1+1\vec{e}_2)\cdot(2\vec{e}_1+1\vec{e}_2)\\
  &=2^2(\vec{e}_1\cdot\vec{e}_1)+2(2)(1)(\vec{e}_1\cdot\vec{e}_2)+1^2(\vec{e}_2\cdot\vec{e}_2)\\
  &=2^2+1=5
\end{align*}

The metric or fundamental tensor allows you to define fundamental properties like lengths and angles in a coordinate space



\chapter{Derivatives of Tensors}



\chapter{Curvature}



\chapter{Bibliography}

%\backmatter
%\addcontentsline{toc}{chapter}{Index}
%\printindex

\end{document}
